\section{Introduction}

A fungal genome submission is judged on its nuclear content. The mitochondrial contig is
assembled because the reads are there, deposited because completeness is a virtue, and
annotated to whatever depth the pipeline happens to reach --- which, for organellar sequence, is
often barely at all. This is defensible: the assembly was not undertaken to describe the
mitochondrion.

It is nonetheless consequential for anyone who later maps RNA to that assembly. Fungal
mitochondria are transcriptionally active, their genomes are intron-rich
\cite{song2020bipolaris}, and in libraries where ribosomal depletion is incomplete they can
account for a substantial share of sequenced material. Reads that map to unannotated
mitochondrial sequence are not merely uncounted; they are reported in categories --- ``no
feature'', ``ambiguous'' --- that most workflows summarise and few interrogate.

Song and colleagues' analysis of the \textit{Bipolaris sorokiniana} mitogenome
\cite{song2020bipolaris} illustrates what a fully annotated fungal mitochondrion contains: a
137,775\,bp genome carrying 12 core protein-coding genes interrupted by 28 introns, 38 tRNAs,
and 52 free-standing open reading frames of unknown function. We asked how much of that is
recoverable for \textit{Pleurotus ostreatus} from RNA-seq alone, using libraries generated for
an entirely different purpose.

\section{Results}

\subsection{The deposited mitogenome is almost unannotated}

The mitochondrial contig CM148777.1 of the BOM\_ss5 assembly is 71,949\,bp with 26.6\% GC. Its
GenBank annotation carries \textbf{13 gene features: 12 tRNAs and one protein-coding gene}
(Table~\ref{tab:mitoann}). A fungal mitogenome typically encodes fourteen
respiratory-chain and ATP-synthase proteins together with the small- and large-subunit
ribosomal RNAs and a variable tRNA set \cite{sandor2018fungalmito}; the closest well-described
agaric, \textit{Agaricus bisporus}, carries fifteen protein-coding genes
\cite{ferandon2013agaricus}. Essentially the entire protein-coding and ribosomal complement is
therefore missing here.

\subsection{A corrected annotation from four independent lines of evidence}

We first established the genetic code empirically rather than assuming it. EMBOSS
\texttt{getorf} yields 51 ORFs of at least 100 codons under translation table 4 against 34
under table 1, confirming that TGA encodes tryptophan --- the mould mitochondrial code.

Combining tRNAscan-SE (25 tRNAs, corroborated by aragorn at 22), barrnap for ribosomal RNA, and
ORF prediction with DIAMOND identification against Swiss-Prot gives \textbf{53 features: 20
protein-coding, 25 tRNA, 2 rRNA and 6 introns} (Fig.~\ref{fig:map}, Table~\ref{tab:mitoann}).
Identified products include cytochrome b, cytochrome \textit{c} oxidase subunits 2 and 3, NADH
dehydrogenase subunits 2, 4 and 5, ATP synthase subunit a, and the small ribosomal subunit
protein uS3m. Two products are intron-encoded endonucleases of the aI5 and aI8 families, which
is itself evidence of intron content: these enzymes are encoded \textit{within} group I introns.

Every feature in the released GFF carries the evidence that produced it, so a reader who
distrusts any one line can discount it without re-deriving the rest.

\subsection{Read assignment improves 2{,}135-fold}

Across all 16 libraries, 21,180,726 alignments fall on the mitochondrion. The stock annotation
assigns \textbf{343} of them. The corrected annotation assigns \textbf{732,325}, a 2,135-fold
increase, and \textbf{40 of 47 quantifiable features are detectably expressed}
(Fig.~\ref{fig:expr}a).

The corrected figure is still only 3.5\% of mitochondrial reads, and the reason is instructive:
barrnap identifies the conserved cores of the ribosomal RNA genes, a few hundred base pairs
each, whereas the transcribed ribosomal regions span several kilobases. Annotation completeness
is not binary, and the residue here is ribosomal rather than protein-coding.

\subsection{Six introns, confirmed by splicing}

Split-read junctions were pooled across all 16 libraries and filtered to those seen in at least
three libraries with at least ten supporting reads, leaving 56 recurrent junctions of 459 raw.

Six of these splice at greater than 50\% efficiency and we regard them as genuine introns
(Table~\ref{tab:introns}, Fig.~\ref{fig:introns}). They total 4,332\,bp, 6.0\% of the
mitogenome. The largest, at 6,796--7,232, is supported by 4,782 reads, is present in
\textbf{all 16 libraries}, and is excised at \textbf{complete efficiency in every one}. Two of
the six coincide with Rfam covariance-model predictions for group I introns, one at
$E=2.7\times10^{-15}$.

Covariance-model search independently identifies nine group I intron candidates on the
mitogenome and no significant group II candidate. This matters for interpretation: group II
introns encode reverse transcriptases ancestral to telomerase and to non-LTR retrotransposon
RTs, whereas group I introns encode homing endonucleases and carry no such connection. The
intron-encoded endonucleases recovered above are consistent with an exclusively group I
complement, and no inference about retroelement or telomerase biology is available from these
data.

\subsection{Most recurrent junctions are artefacts, and they are not randomly placed}

The remaining 50 junctions splice at under 5\% efficiency: for each, the overwhelming majority
of reads crossing the donor site do so without a gap. \textbf{Forty-six of those 50 lie inside
the mitochondrial ribosomal RNA block} (Fig.~\ref{fig:introns}b), where coverage reaches
$10^6$ and a small per-read misalignment rate produces a large absolute number of spurious
gapped alignments.

This has a practical consequence we did not anticipate. Raw split-read counts on the
mitochondrion vary 20--34-fold \textit{between replicates of the same tissue}, which invites
the interpretation that splicing is condition-dependent. It is not: that variance is carried
almost entirely by the artefactual junctions, while the genuine introns splice to completion
uniformly. \textbf{Splicing efficiency, not split-read count, is the quantity that behaves.}
We accordingly make no claim about differential splicing between tissues, and would caution
against such claims from comparable data without an efficiency denominator.

\subsection{Expressed ORFs without characterised homologues}

Twenty of the 51 predicted ORFs have no DIAMOND hit to Swiss-Prot and do not overlap an
annotated feature. Six are detected at five or more reads in at least half the libraries; one
lies inside a ribosomal RNA block and is discarded because its count reflects ribosomal rather
than ORF transcription, leaving \textbf{five credible expressed candidates}
(Table~\ref{tab:orfs}, Fig.~\ref{fig:expr}b). They fall at two loci, near 32.8\,kb and
65.2\,kb, with overlapping predictions at each, so the number of distinct transcribed units is
probably two rather than five.

Absence of a Swiss-Prot hit is weak evidence of novelty --- Swiss-Prot covers 41\% of this
organism's nuclear proteome --- so these are candidates for follow-up, not discoveries.

\section{Discussion}

Three points generalise.

\textbf{Organellar annotation is a systematic blind spot with a measurable cost.} The gap here
is not subtle: one annotated protein-coding gene where there are twenty, and a 2,135-fold
difference in what can be counted. Any study mapping RNA to a fungal assembly inherits this,
and inherits it silently, because unannotated sequence produces no error.

\textbf{RNA-seq is sufficient to repair it.} Nothing here required mitochondrial enrichment,
long reads, or a dedicated experiment. The libraries were prepared for a tissue comparison and
the mitochondrial signal was, from that perspective, contamination. Where poly(A) selection is
incomplete --- common enough --- the resulting non-polyadenylated material is exactly what makes
organellar transcript annotation possible.

\textbf{Split-read counts are not splicing measurements.} The distinction between a junction
and an efficiently spliced junction is the difference between 56 apparent introns and six real
ones, and between an apparent 20--34-fold biological effect and an artefact of coverage. We
would not have separated them without computing the unspliced denominator.

\subsection{Limitations}

The annotation is computational throughout and no gene model has been verified by RACE,
proteomics or targeted sequencing. Protein identifications rest on homology, and the ORF
predictions inherit whatever errors are present in the underlying assembly, which is
contig-level. Coverage is uneven --- median depth 2 outside the ribosomal regions --- so absence
of expression is weak evidence for any individual feature. The intron catalogue is limited to
introns that are spliced in these libraries; a constitutively retained intron is invisible to
this approach by construction. Rfam group I models are sensitive but not specific at the
thresholds used, and seven of the nine covariance-model hits lack independent split-read
support.

\section{Methods}

Culture, RNA extraction, sequencing and read processing are described in the companion resource
paper. In brief, 16 libraries were sequenced as poly(A)-selected 3$'$-tag libraries and aligned
with HISAT2 \cite{kim2019hisat2} to \textit{P.\ ostreatus} BOM\_ss5
(\texttt{GCA\_056149245.1}); alignments were reduced to primary records before analysis.

\textbf{Annotation.} Genetic code determined by comparing EMBOSS \texttt{getorf} yield under
tables 1 and 4. tRNAs from tRNAscan-SE \cite{chan2021trnascan} in organellar mode, cross-checked
with aragorn \cite{laslett2004aragorn}; rRNAs from barrnap \cite{seemann2013barrnap};
protein-coding genes from \texttt{getorf} ORFs identified by DIAMOND \cite{buchfink2015diamond}
blastp against Swiss-Prot \cite{uniprot2023}, retaining the longest ORF where predictions
nested. Reconciliation and GFF output by \texttt{scripts/40\_mito\_annotate.py}.

\textbf{Introns.} Junctions extracted from CIGAR \texttt{N} operations, pooled across libraries
and filtered to $\geq$3 libraries and $\geq$10 reads. Splicing efficiency computed per library
as spliced reads over spliced plus donor-spanning unspliced reads, requiring at least five
informative reads. Group I/II predictions from Infernal \cite{nawrocki2013infernal}
\texttt{cmsearch} against Rfam \cite{kalvari2021rfam} models RF00028, RF00029, RF01998, RF02001
and RF02003.

\textbf{Quantification.} featureCounts \cite{liao2014featurecounts} against the corrected
annotation and, for comparison, against the stock GenBank features.

\clearpage
\section*{Figures}

\begin{figure}[htbp]\centering\includegraphics[width=\textwidth]{fig8_mitogenome_map.pdf}
\caption{\textbf{The deposited mitogenome is almost unannotated.} \textbf{a}, Pooled read depth
across CM148777.1. \textbf{b}, Annotated features: the stock GenBank annotation above, the
corrected annotation below.}\label{fig:map}\end{figure}

\begin{figure}[htbp]\centering\includegraphics[width=\textwidth]{fig9_introns.pdf}
\caption{\textbf{Six introns among fifty-six recurrent junctions.} \textbf{a}, Splicing
efficiency against junction length; point size is supporting reads. \textbf{b}, Artefactual
junctions concentrate inside the ribosomal RNA block. \textbf{c}, The six confirmed introns;
asterisks mark Rfam group I support.}\label{fig:introns}\end{figure}

\begin{figure}[htbp]\centering\includegraphics[width=\textwidth]{fig10_mito_expression.pdf}
\caption{\textbf{Expressed mitochondrial features and candidate ORFs.} \textbf{a}, Pooled reads
for the most highly expressed annotated features. \textbf{b}, Uncharacterised ORFs by position;
those inside ribosomal RNA blocks are discounted.}\label{fig:expr}\end{figure}

\clearpage
\section*{Tables}
\begin{table}[htbp]
\caption{Annotation of the \textit{P.\ ostreatus} BOM\_ss5 mitochondrion (CM148777.1, 71{,}949\,bp).}\label{tab:mitoann}
\footnotesize
\begin{tabular}{@{}lrr@{}}
\toprule
Feature class & Stock GenBank & Corrected (this study) \\
\midrule
Protein-coding (CDS) & 1 & 20 \\
Ribosomal RNA & 0 & 2 \\
Transfer RNA & 12 & 25 \\
Introns & 0 & 6 \\
\textbf{Total} & \textbf{13} & \textbf{53} \\
\bottomrule
\end{tabular}
\begin{tablenotes}[flushleft]\footnotesize\item Corrected features derive from tRNAscan-SE, barrnap, EMBOSS \texttt{getorf} under genetic code 4 with DIAMOND identification, and split-read intron evidence.\end{tablenotes}
\end{table}

\begin{table}[htbp]
\caption{Mitochondrial introns confirmed by split-read evidence (splicing efficiency $>$0.5).}\label{tab:introns}
\footnotesize
\begin{tabular}{@{}rrrrrrl@{}}
\toprule
Start & End & Length & Reads & Libraries & Efficiency & Rfam model \\
\midrule
6796 & 7232 & 436 & 4,782 & 16/16 & 1.0 & group I (E=2.7e-15) \\
19401 & 20688 & 1287 & 146 & 8/16 & 1.0 & group I (E=1.2e-13) \\
10389 & 12124 & 1735 & 124 & 7/16 & 1.0 & -- \\
55069 & 55177 & 108 & 32 & 5/16 & 0.603 & -- \\
23774 & 24482 & 708 & 26 & 5/16 & 1.0 & -- \\
3584 & 3642 & 58 & 24 & 6/16 & 1.0 & -- \\
\bottomrule
\end{tabular}
\begin{tablenotes}[flushleft]\footnotesize\item Efficiency is spliced reads divided by spliced plus donor-spanning unspliced reads, averaged over libraries with at least five informative reads.\end{tablenotes}
\end{table}

\begin{table}[htbp]
\caption{Expressed mitochondrial ORFs with no Swiss-Prot homologue, outside the rRNA blocks.}\label{tab:orfs}
\footnotesize
\begin{tabular}{@{}lrrrcrr@{}}
\toprule
ORF & Start & End & Length (nt) & Strand & Reads & Libraries \\
\midrule
CM148777.1\_27 & 65174 & 66319 & 1146 & + & 584 & 12/16 \\
CM148777.1\_14 & 32764 & 33906 & 1143 & + & 524 & 11/16 \\
CM148777.1\_33 & 65782 & 66177 & 396 & - & 376 & 11/16 \\
CM148777.1\_32 & 66224 & 66730 & 507 & - & 254 & 10/16 \\
CM148777.1\_44 & 32880 & 33218 & 339 & - & 218 & 8/16 \\
\bottomrule
\end{tabular}
\begin{tablenotes}[flushleft]\footnotesize\item One further ORF passed the expression filter but lies inside an mt-rRNA block and is excluded: its read count reflects ribosomal RNA rather than ORF transcription.\end{tablenotes}
\end{table}

